\documentclass[11pt,a4paper]{article}

\usepackage[utf8]{inputenc}
\usepackage[T1]{fontenc}
\usepackage[ngerman]{babel}
\usepackage{lmodern}
\usepackage[margin=2.2cm]{geometry}
\usepackage{xcolor}
\usepackage{enumitem}
\usepackage{booktabs}
\usepackage{tabularx}
\usepackage{tikz}
\usepackage{hyperref}

\definecolor{mika}{HTML}{4F81BD}
\definecolor{yannik}{HTML}{F08A5D}
\definecolor{beide}{HTML}{8E44AD}
\definecolor{rahmen}{HTML}{8C8C8C}

\newcommand{\mika}{\textcolor{mika}{\textbf{Mika}}}
\newcommand{\yannik}{\textcolor{yannik}{\textbf{Yannik}}}
\newcommand{\beide}{\textcolor{beide}{\textbf{Beide}}}


\newcommand{\ap}[7]{%
  \item[\textbf{#1}] \textbf{#2}\par
  \begin{tabularx}{\linewidth}{@{}l X@{}}
    Ziel: & #7 \\
    Verantwortlich: & #3 \\
    Aufwand: & #4 PT \\
    Zeitraum: & #5 \\
    Vorgänger: & #6 \\
  \end{tabularx}
  \vspace{0.3em}
}

\title{Blind Spot: Arbeitspakete}
\author{Gruppe: Mika Brückner, Yannik Köllmann\\[0.3em]
  \small Anwendungspraktikum 3D-Rechnersehen / Projekt \glqq Intelligente Systeme\grqq}
\date{November 2026}

\begin{document}
\maketitle

\section*{Projektüberblick}

Blind Spot ist ein Hinderniswarner für Blinde. Eine Kamera auf Brusthöhe erkennt
Hindernisse im Gehbereich, schätzt ihren Abstand monokular und warnt per Stereo-Ton.
Das System ergänzt den Langstock bei Hindernissen auf Brust- und Kopfhöhe.

Ein vortrainiertes Tiefenmodell dient als Basismodell. Darauf aufbauend verbessern wir
die metrische Genauigkeit auf zwei Wegen parallel: durch Fine-Tuning auf passenden Daten
und durch eine Skalierung über Bodenebene und bekannte Kamerahöhe (Bodenskalierung).
Ein iPhone 12 Pro mit LiDAR liefert Referenztiefen.

\begin{center}
\begin{tabularx}{\linewidth}{@{}llX@{}}
  \toprule
  Meilenstein & Datum & Ergebnis \\
  \midrule
  M0 & 04.11.2026 & Planung fertig \\
  M1 & 18.11.2026 & Testset, Auswertungsframework und Offline-Pipeline mit Basismodell \\
  M2 & 09.12.2026 & Echtzeitbetrieb mit Audio-Feedback, Bodenskalierung integriert \\
  M3 & 06.01.2027 & Fine-Tuning abgeschlossen und bewertet \\
  M4 & 20.01.2027 & Prototyp fertig, alle Varianten ausgewertet \\
  Vorstellung & 03.02.--10.02.2027 & Vortrag und Demo \\
  \bottomrule
\end{tabularx}
\end{center}

Aufwand in Personentagen (PT), nicht in Vollzeittagen. Jedes Arbeitspaket hat eine
verantwortliche Person (\mika{} oder \yannik). Die Ergebnissicherung bearbeiten
beide gemeinsam (\beide).


\clearpage
\section{Work Breakdown Structure}

Ebene~1 ist nach Objekt gegliedert, also nach den Komponenten des Systems.
Ebene~2 ist nach Funktion gegliedert, also nach Tätigkeiten.
Die zeitliche Orientierung folgt im Gantt-Diagramm.

\begin{figure}[h]
\centering
\begin{tikzpicture}[
    font=\footnotesize,
    root/.style={draw=mika, thick, fill=mika!12, rounded corners=2pt,
                 font=\normalsize\bfseries, minimum width=3cm, minimum height=0.8cm},
    comp/.style={draw=rahmen, rounded corners=2pt, fill=rahmen!10,
                 font=\footnotesize\bfseries, text width=2.8cm, align=center,
                 minimum height=0.95cm},
    pkg/.style={draw=rahmen, rounded corners=2pt, anchor=west,
                text width=2.5cm, align=left, minimum height=1.2cm, inner sep=3pt},
    line/.style={draw=rahmen, thick},
  ]


  \def\colw{3.35}
  \node[root] (root) at (2*\colw+1.5, 0) {Blind Spot};


  \foreach \i/\name in {0/{1 Datengrundlage}, 1/{2 Gelernte\\Tiefenpipeline},
                        2/{3 Warnsystem}, 3/{4 Evaluation}, 4/{5 Ergebnis-\\sicherung}} {
    \node[comp] (c\i) at (\i*\colw+1.5, -1.6) {\name};
  }


  \draw[line] (root.south) -- ++(0,-0.4) coordinate (fork);
  \draw[line] (0*\colw+1.5, -0.8) -- (4*\colw+1.5, -0.8);
  \foreach \i in {0,...,4} \draw[line] (\i*\colw+1.5, -0.8) -- (c\i.north);

  \newcommand{\pkgs}[2]{%
    \foreach \t [count=\k] in {#2} {
      \node[pkg] (p#1-\k) at (#1*\colw+0.4, -1.6-1.45*\k) {\t};
      \draw[line] (#1*\colw+0.2, -1.6-1.45*\k) -- (p#1-\k.west);
    }
  }
  \pkgs{0}{{1.1 Kamera-\\kalibrierung\\\yannik},{1.2 SANPO\\aufbereiten\\\yannik},
           {1.3 Eigene\\Aufnahmen\\\mika},{1.4 Annotation\\und Splits\\\yannik}}
  \draw[line] (0*\colw+0.2, -2.075) -- (0*\colw+0.2, -1.6-1.45*4);

  \pkgs{1}{{2.1 Auswahl\\Basismodell\\\mika},{2.2 Boden-\\skalierung\\\mika},
           {2.3 Fine-Tuning\\\yannik},{2.4 Kombination\\beider Wege\\\mika}}
  \draw[line] (1*\colw+0.2, -2.075) -- (1*\colw+0.2, -1.6-1.45*4);

  \pkgs{2}{{3.1 Gehkorridor\\und Warnlogik\\\mika},{3.2 Stereo-\\Audio\\\mika},
           {3.3 Echtzeit-\\Integration\\\mika},{3.4 Fehlalarm-\\filter\\\yannik}}
  \draw[line] (2*\colw+0.2, -2.075) -- (2*\colw+0.2, -1.6-1.45*4);

  \pkgs{3}{{4.1 Auswertungs-\\framework\\\yannik},{4.2 Gesamt-\\auswertung\\\yannik},
           {4.3 Live-Test,\\Fehleranalyse\\\mika}}
  \draw[line] (3*\colw+0.2, -2.075) -- (3*\colw+0.2, -1.6-1.45*3);

  \pkgs{4}{{5.1 Doku-\\mentation\\\beide},{5.2 Abschluss-\\bericht\\\beide},
           {5.3 Vortrag\\und Demo\\\beide}}
  \draw[line] (4*\colw+0.2, -2.075) -- (4*\colw+0.2, -1.6-1.45*3);
\end{tikzpicture}
\caption{Work Breakdown Structure: 5 Komponenten, 18 Arbeitspakete}
\end{figure}

\clearpage


\section{Arbeitspakete}

\subsection*{1 Datengrundlage}
\begin{description}[leftmargin=2.2em, labelwidth=2em, itemsep=0.4em]
  \ap{1.1}{Kamerakalibrierung}{\yannik}{1}{02.11.--06.11.2026}{--}
     {Für das iPhone sind die Kameraparameter (Brennweite und Bildmittelpunkt)
      sowie die Höhe der Kamera am Körper bestimmt und im Repo abgelegt.}
  \ap{1.2}{SANPO aufbereiten}{\yannik}{2}{02.11.--13.11.2026}{--}
     {Eine Teilmenge von SANPO (Brustkamera, Synthetic und ausgewählte Real-Sessions)
      ist geladen und über einen Loader mit Tiefe in Metern lesbar.}
  \ap{1.3}{Eigene Aufnahmen}{\mika}{2}{05.11.--15.11.2026}{1.1}
     {Etwa 10 Testwege innen und außen mit definierten Hindernistypen sind mit
      iPhone-LiDAR aufgenommen, dazu getrennte Trainingsaufnahmen.}
  \ap{1.4}{Annotation und Splits}{\yannik}{2}{12.11.--18.11.2026}{1.2, 1.3}
     {Hindernisse im Testset sind markiert und feste Train-, Val- und Test-Splits
      liegen versioniert im Repo.}
\end{description}

\subsection*{2 Gelernte Tiefenpipeline}
\begin{description}[leftmargin=2.2em, labelwidth=2em, itemsep=0.4em]
  \ap{2.1}{Auswahl des Basismodells}{\mika}{2}{09.11.--18.11.2026}{1.2, 4.1}
     {Zwei bis drei metrische Tiefenmodelle sind auf den Testdaten verglichen und das
      Basismodell ist festgelegt.}
      \newpage
  \ap{2.2}{Bodenskalierung}{\mika}{4}{19.11.--09.12.2026}{1.1, 2.1}
     {Die Tiefe eines beliebigen Modells wird pro Bild über Bodenebene (RANSAC) und
      Kamerahöhe in Meter korrigiert, mit Rückfall bei fehlendem Boden.}
  \ap{2.3}{Fine-Tuning}{\yannik}{5}{19.11.--23.12.2026}{1.4, 2.1}
     {Ein auf SANPO und eigenen Trainingsdaten nachtrainiertes Modell liegt vor und
      ist gegen das Basismodell ausgewertet.}
  \ap{2.4}{Kombination von Fine-Tuning und Bodenskalierung}{\mika}{1}{28.12.2026--06.01.2027}{2.2, 2.3}
     {Fine-Tuning und Bodenskalierung sind kombiniert ausgewertet und es ist
      entschieden, ob das nachtrainierte Modell in den Prototyp eingeht.}
\end{description}

\subsection*{3 Warnsystem}
\begin{description}[leftmargin=2.2em, labelwidth=2em, itemsep=0.4em]
  \ap{3.1}{Gehkorridor und Warnlogik}{\mika}{2}{09.11.--18.11.2026}{1.2}
     {Hindernisse im Gehkorridor lösen offline eine Warnung aus, die naive Baseline
      ohne Korridor ist umgesetzt.}
  \ap{3.2}{Stereo-Audio}{\mika}{2}{23.11.--02.12.2026}{3.1}
     {Richtung wird über links und rechts, Nähe über die Piepfrequenz hörbar
      ausgegeben.}
  \ap{3.3}{Echtzeit-Integration}{\mika}{3}{30.11.--09.12.2026}{2.1, 3.1, 3.2}
     {Kamerastream, Tiefenmodell, Warnlogik und Audio laufen live in einer Schleife,
      die Latenz ist gemessen.}
  \ap{3.4}{Fehlalarmfilter}{\yannik}{2}{10.12.--23.12.2026}{2.2, 3.3}
     {Mindestgröße, Mehrbild-Bestätigung und ein 3D-Korridor über die Bodenebene
      reduzieren Fehlalarme.}
\end{description}

\subsection*{4 Evaluation}
\begin{description}[leftmargin=2.2em, labelwidth=2em, itemsep=0.4em]
  \ap{4.1}{Auswertungsframework}{\yannik}{2}{02.11.--13.11.2026}{--}
     {Ein Skript bewertet jede Variante gleich: Tiefenfehler, Recall, Fehlalarme,
      Richtung und Latenz.}
  \ap{4.2}{Gesamtauswertung}{\yannik}{3}{11.01.--20.01.2027}{2.4, 3.4}
     {Alle Varianten sind auf dem Testset verglichen: naive Baseline, Basismodell,
      Basismodell mit Bodenskalierung, nachtrainiertes Modell und nachtrainiertes
      Modell mit Bodenskalierung.}
  \ap{4.3}{Live-Test und Fehleranalyse}{\mika}{2}{14.01.--20.01.2027}{3.3, 3.4}
     {Ein Live-Test mit Augenbinde ist durchgeführt und Fehlerfälle (Kopfhöhe,
      Treppe, Glas, Personen) sind analysiert.}
\end{description}

\subsection*{5 Ergebnissicherung}
\begin{description}[leftmargin=2.2em, labelwidth=2em, itemsep=0.4em]
  \ap{5.1}{Dokumentation}{\beide}{2}{02.11.2026--27.01.2027}{--}
     {Repo, README und Code-Dokumentation sind laufend gepflegt und zum Projektende
      vollständig.}
      \newpage
  \ap{5.2}{Abschlussbericht}{\beide}{4}{18.01.--02.02.2027}{4.2, 4.3}
     {Der Abschlussbericht mit Methoden, Ergebnissen und Fehleranalyse ist fertig.}
  \ap{5.3}{Vortrag und Demo}{\beide}{3}{21.01.--02.02.2027}{4.2}
     {Vortragsfolien und ein Demo-Video als Absicherung für die Live-Vorführung sind
      fertig.}
\end{description}


\section{Arbeitsverteilung}

\begin{center}
\begin{tabular}{@{}llr@{}}
  \toprule
  Person & Arbeitspakete & Aufwand (PT) \\
  \midrule
  \mika{} Brückner & 1.3, 2.1, 2.2, 2.4, 3.1, 3.2, 3.3, 4.3 & 18 \\
  \yannik{} Köllmann & 1.1, 1.2, 1.4, 2.3, 3.4, 4.1, 4.2 & 17 \\
  \beide{} gemeinsam & 5.1, 5.2, 5.3 & 9 \\
  \midrule
  Anteil pro Person & & Mika 22,5 / Yannik 21,5 \\
  \midrule
  Gesamt & & 44 \\
  \bottomrule
\end{tabular}
\end{center}

\end{document}
